\subsection{Strategy 1: direct length-sum obstructions}
\label{subsec:length-guide}

For a closed triangle $T$ and a rectifiable target $E$, put
$L_E(T)=\mathcal H^1(T\cap E)$.  If the original open roles cover $E$,
then their closures satisfy
\begin{equation}
 \mathcal H^1(E)\le L_E(T_C)+\sum_{i=0}^5L_E(T_i).
 \label{eq:length-principle}
\end{equation}
Strategy~1 contradicts this inequality on one of three targets.  Figure
\ref{fig:strategy1-targets} is a visual index; the skeleton estimate is
always conditional on the branch hypotheses and is never asserted globally.

\begin{figure}[H]
\centering
\begin{tikzpicture}[scale=.84]
  \foreach \x/\title in {0/$\partial H$: length $6$,4/$D$: length $6$,8/$S=\partial H\cup D$: length $12$}{
    \begin{scope}[xshift=\x cm]
      \coordinate (a0) at (1,0);
      \coordinate (a1) at (.5,.866);
      \coordinate (a2) at (-.5,.866);
      \coordinate (a3) at (-1,0);
      \coordinate (a4) at (-.5,-.866);
      \coordinate (a5) at (.5,-.866);
      \coordinate (o) at (0,0);
      \draw[hex] (a0)--(a1)--(a2)--(a3)--(a4)--(a5)--cycle;
      \node[panellabel,align=center] at (0,-1.38) {\title};
    \end{scope}
  }
  \draw[demand] (1,0)--(.5,.866)--(-.5,.866)--(-1,0)--(-.5,-.866)--(.5,-.866)--cycle;
  \foreach \p in {(4,0),(4.5,.866),(3.5,.866),(3,0),(3.5,-.866),(4.5,-.866)}{
    \draw[radialdemand] (4,0)--\p;
  }
  \draw[demand] (9,0)--(8.5,.866)--(7.5,.866)--(7,0)--(7.5,-.866)--(8.5,-.866)--cycle;
  \foreach \p in {(9,0),(8.5,.866),(7.5,.866),(7,0),(7.5,-.866),(8.5,-.866)}{
    \draw[radialdemand] (8,0)--\p;
  }
  \node[figlabel,align=center,text=DemandRed] at (8,-1.93)
    {the full-skeleton estimate is used only under its stated branch hypotheses};
\end{tikzpicture}

\caption{Exact targets and their one-dimensional measures.  The coloring is
explanatory.  The full-skeleton panel is conditional, not a claim that seven
unit triangles can never cover $S$.}
\label{fig:strategy1-targets}
\end{figure}

\subsubsection{The human-level mechanism}

There are three ways for the sum in \eqref{eq:length-principle} to fail.

\begin{enumerate}
\item On $\partial H$, center and vertex classifications give a short trace
table.  A Vd1/Vd2 role pays a half-unit boundary cap, while a supercritical
row forces compensating restrictions on neighboring traces.  This closes
the perimeter branches.
\item On $D$, a nonsupercritical Vd0 role contributes at most one, a
supercritical Vd0 role strictly less than one half, and a T3-like role
strictly less than one half.  If $k\ge2$ rows are T3-like and exactly one
row is supercritical, the total is
\[
 L_D(T_C)+\sum_iL_D(T_i)
 <\frac32+\frac12+\frac{k}{2}+(5-k)
 =7-\frac{k}{2}\le6,
\]
contradicting $\mathcal H^1(D)=6$.
\item On $S$, the center, supercritical, positive-support, and no-support
caps combine only under the hypotheses of the conditional skeleton theorem.
They close $N_+\ge2$ in the CE1/CE2 classes and the indicated hybrid CE2
subcases.
\end{enumerate}

The exact neighboring-arm formula, every trace table, and all endpoint
strictness checks are proved in Section~\ref{sec:length}.

\paragraph{Branches closed by the length argument.}
Section~\ref{sec:length} derives the boundary, radial, diagonal, and
conditional-skeleton contradictions and ends with the exact list in
Proposition~\ref{prop:length-branches}.  It closes every row of
Table~\ref{tab:routing} assigned to route $1$.  Within the hybrid $1+2$ row,
it closes the Vd2 neighboring-midpoint and additional-positive-support
subcases; the complementary exact placements are treated by Strategy~2.
